% Part: many-valued-logic
% Chapter: syntax-and-semantics
% Section: sublogics

\documentclass[../../../include/open-logic-section]{subfiles}

\begin{document}

\olfileid{mvl}{syn}{sub}

\olsection{શાસ્ત્રીય તર્કશાસ્ત્ર~$\LogCL$ના ઉપતર્કશાસ્ત્રો તરીકે બહુમૂલ્ય તર્કશાસ્ત્રો}

સામાન્ય બહુમૂલ્ય તર્કશાસ્ત્રોની વ્યાખ્યા એવી શ્રેણિકો વડે થાય છે જેમાં
$\{\True, \False\}$માંથી લેવાયેલા આર્ગ્યુમેન્ટ માટે સત્યવિધેયનું
મૂલ્ય શાસ્ત્રીય સત્યવિધેયના મૂલ્ય સાથે મળે છે. ચોક્કસ કહીએ તો, આ
તર્કશાસ્ત્રોમાં $x \in \{\True, \False\}$ હોય ત્યારે
$\tf{\lnot}[\Log L](x) = \tf{\lnot}[\LogCL](x)$; અને $\star$ એ
$\land$, $\lor$, $\lif$માંથી કોઈ પણ એક હોય અને
$x, y \in \{\True, \False\}$ હોય ત્યારે
$\tf{\star}[\Log L](x,y) = \tf{\star}[\LogCL](x,y)$.
બીજા શબ્દોમાં, $\lnot$, $\land$, $\lor$, $\lif$ માટેનાં સત્યવિધેયો
$\{\True,\False\}$ પર મર્યાદિત કરીએ તો તે ચોક્કસ શાસ્ત્રીય
સત્યવિધેયો છે.

\begin{prop}\ollabel{prop:mvl-cl}
  ધારો કે બહુમૂલ્ય તર્કશાસ્ત્ર~$\Log L$ની ભાષામાં $\lnot$, $\land$,
  $\lor$, $\lif$ સંયોજકો છે, $\True, \False \in V$ છે, અને તેનાં
  સત્યવિધેયો નીચેની શરતો સંતોષે છે:
  \begin{enumerate}
    \item\ollabel{prop:not} $\tf{\lnot}[\Log L](x) = \tf{\lnot}[\LogCL](x)$ જો $x =
    \True$ અથવા $x = \False$;
    \item\ollabel{prop:land} $\tf{\land}[\Log L](x,y) = \tf{\land}[\LogCL](x,y)$,
    \item\ollabel{prop:lor} $\tf{\lor}[\Log L](x,y) = \tf{\lor}[\LogCL](x,y)$,
    \item\ollabel{prop:lif} $\tf{\lif}[\Log L](x,y) = \tf{\lif}[\LogCL](x,y)$,
    જો $x, y \in \{\True, \False\}$ હોય.
  \end{enumerate}
  ત્યારે ઉપરના ચાર સંયોજકો અને વિધાનચલોથી બનેલા કોઈ પણ સૂત્ર માટે,
  તે સૂત્રના દરેક વિધાનચલને શાસ્ત્રીય મૂલ્ય આપતું~$V$માંનું
  કોઈ પણ મૂલ્યાંકન $\pAssign v$, જેના માટે
  $\pAssign v(p) \in \{\True,\False\}$ હોય, તે હેઠળ
  $\pValue v[\Log L](!A) = \pValue v[\LogCL](!A)$.
\end{prop}
\sourcecorrection{OLMV-001}{સ્થિર અંગ્રેજી મૂળ આ સમતા
બહુમૂલ્ય ભાષાના દરેક સૂત્ર માટે કહે છે, છતાં વધારાના
સંયોજકોનાં સત્યવિધેયો પર કોઈ શરત મૂકતું નથી. તે
મૂલ્યાંકનની દ્વિમૂલ્ય શરત પણ એક મુક્ત વિધાનચલ માટે
જ લખે છે. અહીં વિધાનને દર્શાવેલા ચાર સંયોજકોના
સામાન્ય ખંડ અને સૂત્રના દરેક વિધાનચલ સુધી મર્યાદિત
કર્યું છે; શૂન્યસ્થાનીય સંયોજકનો સમાવેશ કરવા તેના
મૂલ્યની અલગ સમાનતા-શરત જોઈએ. દાખલા તરીકે, બીજા
તર્કમાં અસત્ય-અચલને સત્ય મૂલ્ય આપીએ તો ચાર
સંયોજકોની શરતો છતાં આ સમતા તૂટે છે.}

\begin{proof}
  $!A$ પર અનુમાનથી.
  \begin{enumerate}
  \item જો આણ્વિક સૂત્ર $!A \ident p$ હોય, તો
  $\pValue v[\Log L](!A) = \pAssign v(p) = \pValue
  v[\LogCL](!A)$.
  \item જો $!A \ident \lnot B$ હોય, તો
  \begin{align*}
    \pValue v[\Log L](!A) & = \tf{\lnot}[\Log L](\pValue v[\Log L](!B)) & &\text{\olref[val]{defn:pValue} અનુસાર}\\
    & = \tf{\lnot}[\Log L](\pValue v[\LogCL](!B)) && \text{અનુમાનની પૂર્વધારણાથી}\\
    & = \tf{\lnot}[\LogCL](\pValue v[\LogCL](!B))
&&\text{ધારણા \olref{prop:not} અનુસાર,}\\
&&&\text{કારણ કે $\pValue v[\LogCL](!B) \in \{\True, \False\}$,}\\
& = \pValue v[\LogCL](!A) &&\text{\olref[val]{defn:pValue} અનુસાર}.
\end{align*}
\item જો $!A \ident (!B \land !C)$ હોય, તો
\begin{align*}
  \pValue v[\Log L](!A) & = \tf{\land}[\Log L](\pValue v[\Log L](!B), \pValue v[\Log L](!C)) & &\text{\olref[val]{defn:pValue} અનુસાર}\\
  & = \tf{\land}[\Log L](\pValue v[\LogCL](!B),\pValue v[\LogCL](!C)) && \text{અનુમાનની પૂર્વધારણાથી}\\
  & = \tf{\land}[\LogCL](\pValue v[\LogCL](!B),\pValue v[\LogCL](!C))
&&\text{ધારણા \olref{prop:land} અનુસાર,}\\
&&&\text{કારણ કે $\pValue v[\LogCL](!B),\pValue v[\LogCL](!C) \in \{\True, \False\}$,}\\
& = \pValue v[\LogCL](!A) &&\text{\olref[val]{defn:pValue} અનુસાર}.
\end{align*}
\end{enumerate}
$!A \ident (!B \lor !C)$ અને $!A \ident (!B \lif !C)$ના
કિસ્સા સમાન રીતે સાબિત થાય છે.
\end{proof}

\begin{cor}
  જો બહુમૂલ્ય તર્કશાસ્ત્ર \olref{prop:mvl-cl}ની શરતો સંતોષે,
  $\True \in V^+$ અને $\False \notin V^+$ હોય, તો ઉપરના
  સામાન્ય સૂત્ર-ખંડ પર ${\Entails[\Log L]} \subseteq {\Entails[\LogCL]}$;
  એટલે કે, જો $\Gamma \Entails[\Log L] !B$ હોય, તો
  $\Gamma \Entails[\LogCL] !B$ પણ. ખાસ કરીને,
  $\Log L$ની આ ખંડમાં આવેલી દરેક પુનરુક્તિ શાસ્ત્રીય પુનરુક્તિ પણ છે.
\end{cor}
\sourcecorrection{OLMV-002}{સ્થિર અંગ્રેજી મૂળનો
ફલિતતા-સમાવેશ ભાષાના બધા સૂત્રો માટે જણાવાયો છે.
અગાઉના વિધાનથી તેનો પુરાવો માત્ર બંને ભાષાના
સામાન્ય ચાર-સંયોજક ખંડમાં જ મળે છે; અહીં
નિષ્કર્ષનો વ્યાપ એ પ્રમાણે સ્પષ્ટ કર્યો છે.}

\begin{proof}
  આપણે વિરુદ્ધવિધાન સાબિત કરીએ. ધારો કે $\Gamma \Entails/[\LogCL] !B$.
  ત્યારે એવું કોઈ !!{valuation}~$\pAssign v\colon \PVar \to
  \{\True, \False\}$ છે કે દરેક $!A \in \Gamma$ માટે
  $\pValue v[\LogCL](!A) = \True$ અને
  $\pValue v[\LogCL](!B) = \False$. કારણ કે $\True,
  \False \in V$, એ !!{valuation}~$\pAssign v$ એ~$\Log L$ માટેનું
  !!a{valuation} પણ છે. \olref{prop:mvl-cl} અનુસાર, દરેક
  $!A \in \Gamma$ માટે $\pValue v[\Log L](!A) =
  \True$ અને $\pValue v[\Log L](!B) = \False$.
  $\True \in V^+$ અને $\False \notin V^+$ હોવાથી
  $\pAssign v \Entails[\Log L] \Gamma$ અને
  $\pAssign v \Entails/[\Log L] !B$; એટલે કે,
  $\Gamma \Entails/[\Log L] !B$.
\end{proof}
\end{document}
